\documentclass[12pt]{article}
\usepackage[margin=1in]{geometry}
\begin{document}

\begin{center}
    {\Large\bfseries Short-Trajectory Artifact-Reversal Diffusion for CCTA Calcium Deblooming\par}
    \vspace{1.5em}
    {\large Mahmud Wasif Nafee$^{1}$, Muhan Shao$^{2}$, Christopher Wiedeman$^{2}$, Lin Fu$^{2}$, Bruno De Man$^{2}$, and Ge Wang$^{1}$\par}
\end{center}

\vspace{1em}

\noindent
$^{1}$Biomedical Imaging Center, Center for Computational Innovations, Center for Biotechnology and Interdisciplinary Studies, Department of Biomedical Engineering, School of Engineering, Rensselaer Polytechnic Institute, USA\\
$^{2}$GE HealthCare, Technology and Innovation Center, Niskayuna, NY, USA

\vspace{1em}

\noindent\textbf{Corresponding author:}\\
Ge Wang\\
Biomedical Imaging Center, Center for Computational Innovations, Center for Biotechnology and Interdisciplinary Studies, Department of Biomedical Engineering, School of Engineering, Rensselaer Polytechnic Institute, USA\\
Email: wangg6@rpi.edu

\vspace{1em}

\noindent\textbf{Acknowledgements:} This work was supported in part by the National Institutes of Health under Grants EB032716 and HL151561. We thank GE HealthCare Technology and Innovation Center for their support.

\vspace{1em}

\noindent\textbf{Conflict of interest:} The authors declare no conflicts of interest.

\vspace{1em}

\noindent\textbf{Supplementary Materials:} The synthetic CatSim CCTA dataset and clinical CCTA examples used in this study were obtained from a previously published work. Access to these data is governed by the original authors’ data-sharing policy and any applicable institutional or data-use restrictions. Code and trained model weights may be made available by the corresponding author upon reasonable request, subject to institutional approval and applicable restrictions.

\end{document}
